\scanpage{16}
\mspage{g}
\uline{Corollaire} $p_B^*: \widehat A\mathbin{-}\widehat{(A\times B)}$ commute aux $\underline{\operatorname{Hom}}$
(\unsure{$d^{\circ}$} symétrique pour $p_A^*$)

Il suffit de prendre $\deleted{\unsure{M'=e_{\widehat B},M=}}\quad M'=M=e_{\widehat B}$. On trouve
\[
\underline{\operatorname{Hom}}(L'\boxtimes M',L\boxtimes e_{\widehat B})\simeq\underline{\operatorname{Hom}}(L',L)\boxtimes e_{\widehat B}
\]
en particulier, faisant $L'=e_{\widehat A}$ et changeant $M'$ en $M$
\[
\underline{\operatorname{Hom}}(e_{\widehat A}\boxtimes M,L\boxtimes e_{\widehat B})\xleftarrow{\sim}L\boxtimes e_{\widehat B}
\]
--- on retrouve ici la \deleted{\unsure{flèche can.|flèche}} \unsure{la|cette} flèche canonique
\diagram{\begin{tikzcd}\underline{\operatorname{Hom}}(P,Q)&Q\arrow[l]\end{tikzcd}}
valable en tt \unsure{catégories|catégorie}, \deleted{\unsure{munie de produits}}, \deleted{\unsure{fini}}.

Dans la notation $X^M$, si on identifie $M$ à $e_{\widehat A}\boxtimes M=p_A^*(M)$, et de \unsure{même} $L$ à $L\boxtimes e_{\widehat B}\simeq p_B^*(L)$, on trouve donc
\[
L^M=L\qquad\text{pour }\unsure{L\in\operatorname{Ob}\widehat A|L\operatorname{Ob}\widehat A},\ M\in\operatorname{Ob}\widehat B
\]
\mspage{15}
